\documentclass{article}
\usepackage[utf8]{inputenc}
\usepackage[vietnam]{babel}
\usepackage{amsmath}
\usepackage{booktabs}
\usepackage{enumitem}
\usepackage{pifont}

\begin{document}

\section*{TỔNG QUAN BỘ DỮ LIỆU (DATASET OVERVIEW)}

\subsection{1. Thông tin chung}
\begin{itemize}[label=--]
    \item \textbf{Tên bộ dữ liệu:} Twitter Airline Sentiment Dataset (Dữ liệu đánh giá của khách hàng về các hãng hàng không trên nền tảng Twitter).
    \item \textbf{Quy mô dữ liệu:} Bao gồm tổng cộng \textbf{14,640} bản ghi. Mỗi bản ghi đại diện cho một lượt tweet (bài đăng) của khách hàng.
    \item \textbf{Mục tiêu phân tích:} Đo lường sức khỏe thương hiệu (Brand Sentiment), phân loại cảm xúc khách hàng và bóc tách các nguyên nhân cốt lõi dẫn đến sự không hài lòng đối với các dịch vụ hàng không.
\end{itemize}

\subsection{2. Cấu trúc các biến số chính (Key Variables)}
Dựa trên quá trình phân tích, bộ dữ liệu bao gồm các trường thông tin quan trọng sau:
\begin{description}
    \item[\texttt{tweet\_id}:] Mã định danh duy nhất cho mỗi bài đăng.
    \item[\texttt{airline}:] Tên của hãng hàng không được nhắc đến (Biến phân loại).
    \item[\texttt{airline\_sentiment}:] Nhãn cảm xúc, chia làm 3 lớp: \textit{Positive} (Tích cực), \textit{Neutral} (Trung lập), \textit{Negative} (Tiêu cực).
    \item[\texttt{negativereason}:] Phân loại lý do cụ thể nếu tweet mang cảm xúc tiêu cực (ví dụ: Trễ chuyến, Dịch vụ khách hàng...).
    \item[\texttt{text / text\_preprocessed}:] Nội dung văn bản gốc và nội dung đã qua xử lý ngôn ngữ tự nhiên (NLP).
    \item[\texttt{tweet\_created}:] Dữ liệu chuỗi thời gian (Time-series) ghi nhận mốc thời gian đăng bài.
    \item[\texttt{confidence}:] Chỉ số độ tin cậy của việc gán nhãn cảm xúc hoặc lý do tiêu cực.
\end{description}

\subsection{3. Đặc điểm phân phối dữ liệu (Data Distribution)}

\subsubsection*{A. Khối lượng dữ liệu theo Hãng hàng không}
Dữ liệu bao phủ 6 hãng hàng không lớn với sự phân bổ như sau:
\begin{table}[h!]
\centering
\begin{tabular}{lrr}
\toprule
\textbf{Hãng hàng không} & \textbf{Số lượng Tweets} & \textbf{Tỷ trọng} \\ \midrule
United           & 3,822 & Lớn nhất \\
US Airways       & 2,913 & -- \\
American         & 2,759 & -- \\
Southwest        & 2,420 & -- \\
Delta            & 2,222 & -- \\
Virgin America   & 504   & Thấp nhất \\ \bottomrule
\end{tabular}
\end{table}

\subsubsection*{B. Cơ cấu Cảm xúc (Sentiment Breakdown)}
Bộ dữ liệu cho thấy sự mất cân bằng đặc trưng của ngành dịch vụ trên mạng xã hội:
\begin{itemize}
    \item \textcolor{red}{\textbf{Tiêu cực (Negative):}} 9,178 bài đăng (~62.7\%).
    \item \textcolor{orange}{\textbf{Trung lập (Neutral):}} 3,099 bài đăng (~21.2\%).
    \item \textcolor{green}{\textbf{Tích cực (Positive):}} 2,363 bài đăng (~16.1\%).
\end{itemize}

\subsubsection*{C. Các nguyên nhân tiêu cực phổ biến}
Các "điểm đau" (Pain-points) chính của khách hàng bao gồm: \textit{Customer Service Issue} (Vấn đề dịch vụ), \textit{Late Flight} (Trễ chuyến), \textit{Cancelled Flight} (Hủy chuyến) và \textit{Flight Booking Problems} (Lỗi đặt vé).

\subsection{4. Đánh giá tính khả dụng của bộ dữ liệu}
\begin{itemize}
    \item \textbf{Điểm mạnh:} Cấu trúc tốt cho bài toán phân loại (\textit{Text Classification}). Cột \texttt{negativereason} cho phép phân tích sâu nguyên nhân gốc rễ. Dữ liệu thời gian hỗ trợ tốt cho phân tích biến động (\textit{Spike analysis}).
    \item \textbf{Thách thức:} Văn bản Twitter chứa nhiều từ lóng, viết tắt và lỗi chính tả, đòi hỏi quy trình tiền xử lý (\textit{Preprocessing}) nghiêm ngặt trước khi đưa vào các mô hình Machine Learning như Logistic Regression hay Random Forest.
\end{itemize}

\section{Phân tích Nguyên nhân Tiêu cực (Root Cause Analysis)}
Phân tích lý do dẫn đến đánh giá tiêu cực (Negative Reasons) cho thấy các "điểm đau" (Pain points) cốt lõi:

\begin{itemize}
    \item \textbf{Vấn đề chung lớn nhất:} \textit{Dịch vụ khách hàng kém (Customer Service Issue)} là nguyên nhân số 1 gây phẫn nộ ở hầu hết các hãng (American, Southwest, US Airways, United, Virgin America).
    \item \textbf{Các nguyên nhân kỹ thuật:} \textit{Chuyến bay bị trễ (Late Flight)} và \textit{Hủy chuyến (Cancelled Flight)} là các nguyên nhân phổ biến tiếp theo. Riêng đối với hãng \textbf{Delta}, "Late Flight" là nguyên nhân bị phàn nàn nhiều nhất, vượt qua cả Customer Service.
    \item \textbf{Đặc thù hãng bay:} Virgin America dù ít bị phàn nàn nhất, nhưng vấn đề nổi cộm của họ lại nằm ở khâu đặt vé (\textit{Flight Booking Problems}).
\end{itemize}

\section{Phân tích Ngôn ngữ \& Từ khóa (NLP \& N-grams Analysis)}
Mô hình Machine Learning (Logistic Regression \& Random Forest với TF-IDF/CountVectorizer) đạt độ chính xác khoảng \textbf{81.8\%}, bóc tách các từ khóa mang tính biểu tượng:

\begin{itemize}
    \item \textbf{Từ khóa kích hoạt rủi ro (Hate Drivers):} Các từ thể hiện sự thất vọng và chờ đợi như: \texttt{suck}, \texttt{worst}, \texttt{screw}, \texttt{terrible}, \texttt{unacceptable}, \texttt{ruin}, \texttt{delay}, \texttt{hour}, \texttt{stuck}, \texttt{rude}. Đặc biệt từ \texttt{hour} và \texttt{delay} có trọng số rất cao.
    \item \textbf{Đặc trưng theo hãng:}
    \begin{itemize}
        \item \textbf{Southwest \& American:} Nhắc nhiều đến \texttt{bag}, \texttt{luggage}, \texttt{hold} (hành lý và giữ máy).
        \item \textbf{Virgin America:} Liên quan đến \texttt{site}, \texttt{website}, \texttt{seat} (lỗi hệ thống web/chỗ ngồi).
    \end{itemize}
\end{itemize}

\section{Phân tích Chuỗi thời gian (Time-Series \& Spike Detection)}
Sử dụng thuật toán kiểm tra Z-score ($Z > 3$) để tìm kiếm các "đỉnh" (spikes) khủng hoảng truyền thông.

\begin{quote}
    \textbf{Kết quả:} Thuật toán ghi nhận \textit{"No strong z-score spikes found (z>3)"}. 
\end{quote}
Điều này cho thấy sự tiêu cực diễn ra đều đặn và liên tục hàng ngày, thay vì bùng phát do một sự kiện thảm họa cụ thể nào đó.

\hr
\section*{ \ding{110} Kết luận \& Đề xuất hành động}
\begin{enumerate}
    \item \textbf{Cải tổ bộ phận CSKH là ưu tiên số 1:} Do lỗi dịch vụ áp đảo lỗi kỹ thuật, các hãng (đặc biệt là US Airways và American) cần đào tạo lại thái độ nhân viên và tối ưu hóa thời gian phản hồi để chặn đứng các cụm từ như "hold", "hour".
    \item \textbf{Hành động theo đặc thù thương hiệu:}
    \begin{itemize}
        \item \textbf{Virgin America:} Nâng cấp hệ thống Website và luồng đặt vé.
        \item \textbf{Southwest \& American:} Rà soát quy trình quản lý Hành lý (Luggage/Baggage handling).
        \item \textbf{Delta:} Tập trung giải quyết tỷ lệ trễ chuyến (Late Flight).
    \end{itemize}
\end{enumerate}

\end{document}